% =========================================================
% Introduction — Formal benchmark framing, MineNPC-Task (grounded)
% =========================================================
\section{Introduction}

Building truly capable AI companions for open-world games requires more than one-shot instruction following, it requires agents that can \emph{plan}, \emph{clarify}, and \emph{remember} in order to operate effectively in dynamic, long-horizon environments. However, existing evaluation practices fall short: many rely on overly prescriptive prompts or grant agents privileged access to hidden environment state, artificially inflating their apparent competence and undermining fair comparison across models. This highlights the need for a \emph{transparent, reproducible benchmark} that captures the pressures of mixed-initiative interaction~\cite{Lai2022MixedInitiative}, without granting unfair advantages. Recent efforts, such as the Minecraft Universe (MCU) benchmark~\cite{MCU2023,MCU2023arxiv}, move in this direction by introducing scalable, composable tasks in open-world environments and emphasizing repeatable, human-aligned evaluation.


We introduce \textsc{MineNPC-Task}, a practical benchmark for evaluating memory-aware, mixed-initiative LLM agents in \emph{Minecraft}. Rather than synthetic prompts, tasks are elicited from expert co-play: during our formative and evaluation sessions, experienced players issued real requests, which we normalized into compact templates with explicit preconditions and a small set of slot parameters, and paired with simple machine-checkable validators. The benchmark runs inside a Mineflayer envelope so perception and action are limited to public, in-game APIs; a bounded-knowledge policy forbids admin commands, global map introspection, and scans beyond loaded chunks. The aim is straightforward: a clean, reusable setup that others can run, swap in different models, and obtain comparable numbers, while avoiding hidden shortcuts that have complicated evaluation in prior \emph{Minecraft}-based agents \citep{johnson2016the,fan2022minedojo,wang2023voyager}.

To execute tasks reproducibly we use a model-agnostic evaluation framework (Section~\ref{sec:framework}). The agent presents a brief plan preview that breaks the request into a handful of subtasks, asks a targeted clarifying question only when a slot is unbound (for example, a tool variant or a search radius), acts through Mineflayer skills, and is judged solely on in-world evidence drawn from inventory and equipment deltas, position changes, nearby entities and blocks within loaded chunks, and recent chat. This mixed-initiative pattern aligns with recent agent designs that combine planning with lightweight memory and reflection while keeping behavior legible to human partners \citep{10.1145/3586183.3606763, sarch2024vlm, Hou_2024}. Short scenario walkthroughs (Section~\ref{sec:scenario-walkthroughs}) illustrate how planning, clarification, and reuse of prior context work together.

As an initial snapshot, we instantiate the framework with GPT-4o and run live co-play with \textbf{8} experienced players (Section~\ref{sec:evaluation}). Across \textbf{44} user-authored tasks and \textbf{216} subtasks, we observe \textbf{71} subtask failures (approximately \textbf{33\%}) and report where things went wrong, including code execution, inventory handling, referencing, and navigation, alongside common recoveries such as clarifying a slot, simplifying a goal, or constraining location. Participants rated interaction quality and interface usability positively, with mixed views on personalization and completion, consistent with the need for stronger memory scaffolding. We intentionally omit ablations or model comparisons. The intent is to provide a compact, user-authored benchmark that others can extend. For context, our focus complements broader surveys on AI in games and adaptive interaction \citep{yannakakis2018artificial, Riedl2021} and ongoing efforts to evaluate reasoning in live games \citep{hu2024gamearena}.


The following are the contribution of this work:
\begin{itemize}[leftmargin=*, itemsep=2pt, topsep=2pt]
  \item \textbf{\textsc{MineNPC-Task} suite:} a user-authored benchmark for \emph{Minecraft}, normalized into templates with explicit preconditions and paired with simple, machine-checkable validators under a bounded-knowledge policy. See Appendix~\ref{app:tasksuite} and Section~\ref{sec:setting}. Compared with prior \emph{Minecraft} agents and datasets, we emphasize human-elicited goals, public-API constraints, and validator-backed judging \citep{johnson2016the,fan2022minedojo,wang2023voyager}.
  \item \textbf{Evaluation framework:} a model-agnostic procedure that enforces plan previews and single-turn clarifications when required, constrains perception and action to Mineflayer APIs, and produces validator-backed outcomes suitable for controlled comparisons across LLMs. See Section~\ref{sec:framework}. The design follows calls for transparent, reproducible evaluation in interactive systems \citep{Jennings2024,hu2024gamearena}.
  \item \textbf{Empirical snapshot with GPT-4o:} co-play results from \textbf{8} experts over \textbf{44} tasks, comprising \textbf{216} subtasks with \textbf{71} failures (approximately \textbf{33\%}), along with observations about where mixed-initiative interaction and lightweight memory help—and where brittleness remains. See Section~\ref{sec:evaluation}.
\end{itemize}